\section{Related Work}
\label{sec:related}

We position this SoK against three prior literature clusters: surveys
and taxonomies of zero-knowledge proof systems, prior work on
privacy-preserving group protocols, and prior SoKs in adjacent
privacy-preserving primitive areas. The first cluster overlaps
directly with our taxonomic ambition; the second and third overlap
with the problem class.

\subsection{Surveys and taxonomies of zero-knowledge proof systems}

The cryptographic literature contains substantial survey work on
zero-knowledge proof systems. The Nitulescu introduction to
zk-SNARKs~\cite{nitulescu-2020-snarks} indexes by construction
family (Groth16, PLONK, STARKs) and by setup posture
(per-circuit / universal / transparent), and is the canonical
introductory taxonomy. Sun et al.'s survey of zero-knowledge proofs
in blockchain settings~\cite{sun-2021-zkp-blockchain} indexes by
proof-system primitive and discusses deployment context as a
secondary lens.
The Bulletproofs lineage~\cite{bunz-2018-bulletproofs} marks the IPA
side of the construction taxonomy; later transparent-setup IPA
variants inherit its structure.

What distinguishes the present work from these surveys is the
indexing axis. Existing surveys treat the proof system as the
primary object of interest and the deployment context as a
consequence; we invert this and treat the deployment constraint
(anchor chain, host functions, fee economics) as the primary axis,
with the proof system following as a derived choice. This inversion
is what makes the cross-axis dependency analysis of
\S\ref{sec:dimensions-deps} legible and what justifies the
profile-conditional recommendations of
\S\ref{sec:recommendations}. A purely cryptographic taxonomy would
not produce \S\ref{sec:rec-jurisdiction}'s ``none of C1--C6''
non-recommendation, because the cryptographic-taxonomy lens does not
weight operational constraints.

\subsection{Privacy-preserving group protocols}

The closest prior work to the registry abstraction of
\S\ref{sec:problem} is Signal's Private Group
System~\cite{chase-2020-signal-pgs}, which closes group-membership
metadata against an honest-but-curious application-layer operator
using anonymous credentials. The Signal construction differs from
the registry abstraction in two respects. First, it places an
application-layer operator in the trust path; the registry
abstraction's contribution is precisely to remove that operator.
Second, it does not anchor state on a public ledger; this is a
property the operator-removal trade is in exchange for. The two
constructions are complementary rather than competing: an MLS-shaped
deployment~\cite{rfc9420-mls} using Signal's credential machinery
for membership and a registry of the kind surveyed here for state
anchoring would combine their guarantees.

The Coconut threshold-credential
construction~\cite{sonnino-2019-coconut} sits in an adjacent slot:
it provides verifiable disclosure of attributes via threshold-issued
credentials, with operator-mediated issuance that the registry
abstraction does not require. Coconut and the registry abstraction
are composable in deployment patterns where credential issuance and
state advancement are separate concerns.

Tornado Cash~\cite{pertsev-2019-tornado} and
Semaphore~\cite{semaphore-2024} are direct instantiations of the
registry abstraction with single-relation, fixed-policy
configurations. They sit at C2 in the taxonomy and have driven the
substantial ecosystem of BN254-precompile-aware tooling that
configurations C2 and C5 inherit. Both pre-date the host-function
heterogeneity that makes this paper's taxonomy necessary.

\subsection{SoKs in adjacent primitive areas}

The SoK genre has been applied to several adjacent privacy-preserving
primitive classes: privacy-preserving payments and mixers,
anonymous credentials, decentralized identity, and mixnets. These
prior SoKs have established the methodology of comparing
constructions against an explicit threat model with profile-conditional
recommendations, which the present work follows. The methodological
debt is direct.

The substantive distinction from each:
\begin{itemize}
  \item Privacy-preserving-payments SoKs ---
    Bonneau et al.~\cite{bonneau-2015-sok-bitcoin} as the foundational
    cryptocurrency-research SoK, and the subsequent
    confidential-transaction and mixer literature it
    cites --- treat a different primitive: value transfer with
    confidential amounts and unlinkable senders. Tornado Cash sits in
    both that primitive class and ours; the per-payment relation
    differs structurally from the per-state-transition relation we
    treat as canonical.
  \item Anonymous-credentials work in the
    Camenisch--Lysyanskaya~\cite{camenisch-2001-cl} tradition
    indexes along a different axis: issuer / prover / verifier
    roles, predicate expressivity, and revocation handling. The
    registry abstraction does not have an issuer in the credential
    sense; the Coconut~\cite{sonnino-2019-coconut} threshold variant
    is the closest credential construction to our setting and, as
    noted above, composes with rather than competes with the
    registry.
  \item Mixnet SoKs ---
    Sampigethaya and Poovendran's survey~\cite{sampigethaya-2006-mix}
    is the canonical reference --- target a different threat model:
    network-observer-only adversaries with active routing
    capability. The registry abstraction's threat model includes the
    chain-validator as a party, which mixnet literature
    typically excludes by construction.
  \item Decentralized-identity surveys ---
    M\"uhle et al.~\cite{muhle-2018-ssi} catalogue the
    self-sovereign-identity component space --- assemble the
    credential, registry, and revocation primitives into composite
    systems. The present work treats one of those primitives (the
    registry) in isolation.
\end{itemize}

\subsection{Position}

Our contribution is the deployment-constraint-indexed taxonomy of
\S\ref{sec:dimensions}, the cross-axis dependency analysis of
\S\ref{sec:dimensions-deps}, and the profile-conditional
recommendations of \S\ref{sec:recommendations}. Prior work has either
(a) surveyed the cryptography without the deployment constraints, or
(b) constructed individual systems that instantiate one point in our
design space. We are not aware of a prior SoK that takes the
metadata-hiding-group-registry primitive as the unit of comparison
and indexes by deployment constraint. The present work occupies that
position.
